\documentclass[10pt,twocolumn,letterpaper]{article}

\usepackage{cvpr}
\usepackage{times}
\usepackage{epsfig}
\usepackage{graphicx}
\usepackage{amsmath}
\usepackage{amssymb}
\usepackage{dblfloatfix}
\usepackage{placeins}
\usepackage{fancyhdr}

\usepackage[breaklinks=true,bookmarks=false]{hyperref}

\cvprfinalcopy

\def\cvprPaperID{****}
\def\httilde{\mbox{\tt\raisebox{-.5ex}{\symbol{126}}}}

\ifcvprfinal
\pagestyle{fancy}
\fancyhf{}
\fancyfoot[C]{\thepage}
\renewcommand{\headrulewidth}{0pt}
\renewcommand{\footrulewidth}{0pt}
\fi

\begin{document}

\title{Parallel Logic Expert with Operator Routing}

\author{
Steven Cleasby-Mayeda\\
Oregon State University\\
{\tt\small cleasbys@oregonstate.edu}
{\tt\small steven.cleasby.mayeda@gmail.com}
}
\date{AI535 Deep Learning \\ 27 September 2026}

\maketitle
\begin{abstract}
We introduce a logic-augmented transformer pathway that routes token-level
representations into fixed fuzzy logical operators (AND/OR/NOT) or straight-through-estimation,
so the network learns compositional structure rather than an ensemble of expert networks
with shared architecture design. At each backbone layer, the hidden states are
projected to a lower dimensional K, Q, which the running logic state queries
(cross-attention). The attention matrix is routed, and local
windows over either token or gate-channel axes are composed.
The final logic state is fused with the backbone features before classification.
The proposed architecture augmentation imposes structured, interpretable compositions over
token features. A key design dimension is the axis of composition:
intra-token compositions (feature axis) compares features within a token, while
inter-token compositions (sequence axis) compares features across neighboring
tokens. We evaluate these design choices with Llama-3.1-8B under matched
parameter budgets over three seeds. On ProofWriter with the rulebase in the
input, adding a parallel branch with a fusion head to the backbone gives a
consistent gain: a No-Gate Control (an MLP branch with the same fusion head, no
logic operators or cross-attention) beats a linear-head Baseline by
$+1.4$ pp in every seed (SE 0.2 pp). The logic operators add nothing on top of
that: no routed variant beats the No-Gate Control (inter-token $G=8$:
$-1.3$ pp, intra-token: $-0.4$ pp, and $-0.2$ pp for routed without cross-attention, which
matches the control's used parameters within 0.05\%), and removing
cross-attention changes nothing measurable. Routing entropy stays at 97--98\% of its maximum and a
learned fusion weight shrinks, so the model largely routes uniformly and
down-weights the branch. On MNLI, routed and No-Gate variants are also
indistinguishable. The gain comes from the fusion-head path, not from operator
routing, at this scale.
\end{abstract}

\section{Introduction}

\begin{figure*}[tbp]
\centering
\includegraphics[width=0.75\textwidth]{plots/arch1.png}
\caption{Our architecture, backbone hidden-states feed to logic expert, and fused back with
the backbone features before classification}
\label{fig:architecture}
\end{figure*}

Transformers provide strong performance on a wide variety of tasks \cite{vaswani2023attentionneed},
both chain-of-thought (CoT) prompting and reasoning models can greatly
improve performance on reasoning tasks
\cite{wei2023chainofthoughtpromptingelicitsreasoning,Guo_2025}. However,
many CoT use N-shot prompting. Additionally, both CoT and reasoning models
utilize reasoning by iterative, and extended inference, which requires far more
tokens. This increases inference costs for many tasks
\cite{wei2023chainofthoughtpromptingelicitsreasoning}.
Our approach instead introduces a novel architecture for
structured computation within the latent space, no additional
tokens needed. We will test how this constrained, logic-inspired
pathway affects performance on reasoning problems and generalization.

Mixture-of-Experts (MoE) models
\cite{6797059,shazeer2017outrageouslylargeneuralnetworks,fedus2022switchtransformersscalingtrillion}
introduce conditional computation by routing tokens to specialized subnetworks, or weighting
the results of several/all experts.
Our method introduce conditional computation by routing, weighting
tokens to logical operators. In this formulation,
the operator routing is learned rather than expert weights.
When tied with attention, this method may allow robust routing
in the logic space. This space has forced sparsity due to softmax
and local compositions (local intra-token or inter-token windows),
which may allow emergence of specialized operator-feature connections.

Our activation metrics allows learned logic compositions to be interpreted.
Routing entropy and fusion $\alpha$ parameters allow us to measure the
utilization of the logic expert. Experiments on logical
reasoning benchmarks compare the Baseline, the No-Gate Control network
(a parallel MLP branch with matched parameter count for fair comparison), and the
logic expert variants.
A core design question is how the local operator windows should be applied in the
gate compositions. Our first implementation applies capacity on the channel axis,
(neighboring channels on one token). The next design applies
capacity on the token axis, each channel compares the same indexed feature on
neighboring tokens (channels remain independent). This axis change could
affect which feature relationships are learnable, and we test whether it
affects performance.
Figure~\ref{fig:architecture} provides a high-level view of the
architecture.

This paper makes three contributions. (1) Inter-token compositional routing with
operator reductions and residual logic-state updates. (2) Controlled, seed-paired
ablations over gate count, window axis, fusion $\alpha$, and cross-attention
against a parameter-matched control. (3) An attribution result: on ProofWriter
the parallel branch with its fusion head beats a linear-head Baseline by
$+1.4$ pp in every seed, but on ProofWriter and MNLI operator routing does not
improve over the parameter-matched control, and
mechanism metrics (near-maximal routing entropy, a shrinking learned $\alpha$)
indicate the branch is barely specialized.

\section{Related Work}

Logical and logic-inspired neural architectures span a wide variety of designs,
from classic weightless systems to modern differentiable operator networks
\cite{Aleksander2009ABI}.
Some architectures utilize explicit logical, or logic-inspired
differentiable operators \cite{petersen2022deepdifferentiablelogicgate} as an
alternative to dense transformations, and they provide evidence that structured
logical computation can be more efficient and interpretable than typical
transformers. Our approach builds on this motivation, but integrates operator routing
directly into transformer hidden-state rather than standalone units.

Deep Differentiable Logic Gate Networks combine real-valued logical operators
with continuous relaxations (differentiable approximations of discrete
logical gates used in gradient descent) and can achieve very efficient inference
(one million+ MNIST images/s on a single CPU core)
\cite{petersen2022deepdifferentiablelogicgate}.
Logical Neural Networks (LNNs) \cite{riegel2020logicalneuralnetworks}
map neural units to components of weighted logical formulas, yielding highly
interpretable representations and supporting logic-consistent reasoning.
LNNs use explicitly logic-structured model semantics.
Differentiable Logic Machines (DLMs) \cite{zimmer2023differentiablelogicmachines}
define a continuous relaxation of first-order logic program spaces (via
predicate-weighted parameterization) and learn interpretable logic programs
with gradient descent. In contrast to these approaches, which operate over
structured relational inputs or operations on the input space,
our method integrates logical composition in the feature space,
it learns logical compositions over token features.

MoE architectures use a learned gating mechanism to route
inputs to specialized expert subnetworks. In adaptive mixtures of local experts,
the objective for each expert is decoupled from other experts, enabling
specialization over different input regions, and thus allow both specialized and
general networks \cite{6797059}. These models can also improve inference efficiency
via forced sparse activations.
Conditional computation methods such as MoE and sparse
routing motivates learned token-dependent routing decisions
\cite{shazeer2017outrageouslylargeneuralnetworks,fedus2022switchtransformersscalingtrillion}.
Our model is closest to this line of work, but uses compact differentiable
operators instead of expert-token dispatch among independent expert sub-networks.

Fuzzy-set theory provides continuous, differentiable functions \cite{ZADEH1965338},
We use fuzzy-set operations to estimate logical operations,
which allows training with any off-the-shelf gradient-based optimizer.
Our architecture is contrasted by Neural-Symbolic networks,
where explicit logic is composed in the non-compressed input space
(for example, decision trees or markov-networks)
\cite{besold2017neuralsymboliclearningreasoningsurvey}.

\section{Methods}

\subsection{Architecture and State Representation}

The logic-augmented architecture introduces the logic-expert pathway parallel to
the transformer backbone. At each layer $t$, the backbone produces hidden states:
\begin{equation}
H_t \in \mathbb{R}^{B \times S \times H}.
\end{equation}
In parallel, the logic expert maintains logic states which are updated at every layer:
\begin{equation}
L_t \in \mathbb{R}^{B \times S \times L},
\end{equation}

At each layer, the logic state attends to the backbone via multi-head cross-attention,
producing attended logic features:
\begin{equation}
\tilde{L}_t.
\end{equation}

After the final layer, the logic state is fused with the backbone:
\begin{equation}
H_{\text{fused}} = H_{\text{last}} + \alpha W_f L_{\text{last}},
\label{eq:fusion}
\end{equation}
where $\alpha$ controls the contribution of the logic pathway.

\paragraph{Notation.}
$\pi_t$: routing weights, $s_t$: gate activations, $u_t$: routed gate activations,
$B$: batch size, $S$: sequence length, $H$: backbone hidden dimension,
$L$: logic-state dimension, $G$: gates per operator family (AND/OR/NOT), with
total routed channels $3G$, $C_t$: local window size at layer $t$,
$\tau$: routing softmax temperature.

\subsection{Cross-Attention into Logic Space}

At each layer, logic tokens act as queries, while backbone hidden states are projected
to keys and values. This produces attended logic features:
\begin{equation}
\tilde{L}_t := \text{CrossAttention}(L_{t-1}, H_t).
\end{equation}

These features serve as the input to the routing and composition stage.

\subsection{Routing and Gate Computation}

Routing and gate activation use separate learned projection matrices: $W_r$
produces routing logits, while $W_g$ produces gate activations.

Routings are computed from $\tilde{L}_t$:
\begin{equation}
R_t := \tilde{L}_t W_r, \quad W_r \in \mathbb{R}^{L \times 3G}.
\end{equation}

Routing weights are:
\begin{equation}
\pi_t[b,s,j] = \frac{\exp(R_t[b,s,j]/\tau)}{\sum_{k=1}^{3G} \exp(R_t[b,s,k]/\tau)}.
\end{equation}

Gate activations are computed as:
\begin{equation}
s_t := \sigma(\tilde{L}_t W_g), \quad W_g \in \mathbb{R}^{L \times 3G}.
\end{equation}

The routed soft activations are:
\begin{equation}
u_t^{\mathrm{soft}} := \pi_t \odot s_t.
\end{equation}
In soft mode, $u_t := u_t^{\mathrm{soft}}$. In STE mode (forward binarized), we use
\begin{equation}
u_t^{\mathrm{hard}} := \mathbf{1}[u_t^{\mathrm{soft}} \ge 0.5], \quad
u_t := u_t^{\mathrm{soft}} + \operatorname{sg}(u_t^{\mathrm{hard}} - u_t^{\mathrm{soft}}).
\end{equation}

We split:
\begin{equation}
[u_t^{\text{AND}}, u_t^{\text{OR}}, u_t^{\text{NOT}}] := \operatorname{split}_{G,G,G}(u_t).
\end{equation}

\subsection{Windowed Operator Composition}

Let $v \in [0,1]^{B \times S \times G}$ denote one operator group.

Windows are causal and defined along either:
inter-token axis (sequence), or intra-token axis (channels).

For inter-token composition:
\begin{equation}
v_{\text{win}} \in \mathbb{R}^{B \times S \times G \times C_t}.
\end{equation}

Operators are:
\begin{align}
\text{AND}(x) &= \prod_{i=1}^{C_t} x_i, \\
\text{OR}(x) &= 1 - \prod_{i=1}^{C_t}(1 - x_i), \\
\text{NOT}(x) &= 1 - x.
\end{align}

Applying to routed groups:
\begin{align}
a_{b,s,g} &:= \text{AND}(\operatorname{win}(u_t^{\text{AND}})_{b,s,g,:}), \\
o_{b,s,g} &:= \text{OR}(\operatorname{win}(u_t^{\text{OR}})_{b,s,g,:}), \\
n_{b,s,g,:} &:= \text{NOT}(\operatorname{win}(u_t^{\text{NOT}})_{b,s,g,:}).
\end{align}

AND and OR reduce to $[B,S,G]$.

NOT is applied elementwise and flattened to $[B,S,G \cdot C_t]$, to preserve
negation operation structure.

The concatenated representation has dimension:
\begin{equation}
G + G + G \cdot C_t = G(C_t + 2).
\end{equation}

\subsection{Layer Update and Residual Flow}

The concatenated operator outputs are projected:
\begin{equation}
\Delta L_t := \text{Linear}(\text{concat}(a, o, n)).
\end{equation}

The logic state is updated via:
\begin{equation}
L_t = \mathrm{RMSNorm}(L_{t-1} + \Delta L_t).
\end{equation}

\subsection{Output Fusion}

After the final layer, fusion is applied as in Eq.~\ref{eq:fusion}.
$\alpha$ may be fixed or learned, it directly controls
the contribution of the logic expert.

\section{Experiments}

\begin{table*}[tbp]
\centering
\normalsize
{\renewcommand{\arraystretch}{1.2}
\setlength{\tabcolsep}{6pt}
\resizebox{\textwidth}{!}{%
\begin{tabular}{|l|c|c|c|c|c|}
\hline
Model & Val Accuracy & Val Loss & Fusion $\alpha$ (final) & Norm.\ Routing Entropy & Grad Norm \\
\hline\hline
Baseline & 0.925 $\pm$ 0.005 & 0.198 $\pm$ 0.003 & -- & -- & 28.5 \\
No-Gate Control & \textbf{0.939 $\pm$ 0.004} & \textbf{0.167 $\pm$ 0.014} & 0.01 (fixed) & -- & 20.3 \\
Routed (intra, G=8) & 0.936 $\pm$ 0.004 & 0.194 $\pm$ 0.023 & 0.01 (fixed) & 0.972 & 19.6 \\
Routed (inter, G=8) & 0.927 $\pm$ 0.012 & 0.203 $\pm$ 0.021 & 0.01 (fixed) & 0.972 & 22.9 \\
Routed (inter, G=8) & 0.936 $\pm$ 0.005 & 0.179 $\pm$ 0.021 & 0.1 (fixed) & 0.972 & 22.2 \\
Routed (inter, G=8) & 0.938 $\pm$ 0.011 & 0.172 $\pm$ 0.007 & 0.0077 (learned from 0.01) & 0.973 & 20.9 \\
Routed (inter, G=16) & 0.934 $\pm$ 0.002 & 0.198 $\pm$ 0.022 & 0.01 (fixed) & 0.976 & 20.3 \\
Routed (inter, G=8), no cross-attn & 0.937 $\pm$ 0.010 & 0.178 $\pm$ 0.031 & 0.01 (fixed) & 0.974 & 21.7 \\
\hline
\end{tabular}
}}
\caption{ProofWriter validation results with the rulebase in the input, mean
$\pm$ SD over 3 seeds (5,000 validation examples per seed). Normalized routing
entropy is $H/\ln(3G)$ (1 = uniform routing). Grad norm is the mean over the
final epoch. Abbreviations: Baseline = backbone with a linear head on the last
token (no parallel branch). No-Gate Control = parallel MLP branch without
cross-attention, operator routing, or compositional gating, with the same fusion
and normalized head as the routed models (its used parameters match the routed
model without cross-attention, Table~\ref{tab:param_counts}). Routed = logic
expert variants.}
\label{tab:main_results}
\end{table*}

\subsection{Setup}
We use Llama-3.1-8B \cite{grattafiori2024llama3herdmodels}
as the shared backbone for all variants.
We compare: (i) a Baseline
model, (ii) a No-Gate Control (parameter-matched), and
(iii+) the logic expert variants. The No-Gate Control serves as a
parameter-matched control to separate routing effects from pure
parameter-count gains.
All variants use the 3 label classifier head.
The original Llama head is not used.
In this control, the parallel branch excludes cross-attention, operator
routing, and compositional gating. Its MLP width is set to match the parameter
count of the routed layers. The cross-attention modules are still constructed
in the control but never called, so its trainable-parameter count includes
71M unused weights. The parameters it actually uses match the routed model
without cross-attention (Table~\ref{tab:param_counts}). All experiments use the AdamW
\cite{loshchilov2019decoupledweightdecayregularization}
optimizer and cross-entropy loss, with full fine-tuning of the backbone in
bfloat16. Every configuration is run with three seeds (42, 777, 123). The seed
selects the training and validation subsets as well as the initialization, so
comparisons between models are seed-paired. We call a difference supported only
if its seed-paired mean is positive and exceeds two standard errors. The main
hypothesis (routed inter-token $G=8$ beats the No-Gate Control on ProofWriter)
and this test were fixed before the runs. Each run used one NVIDIA H100 80GB
or H200 141GB GPU on the Oregon State University HPC cluster. Code, per-run
logs and the scripts that generate every table and figure are in the public
repository (\url{https://github.com/Ancorro/logic-expert-public}).

\subsection{Logical Reasoning -- ProofWriter}
We use the ProofWriter \cite{tafjord2021proofwritergeneratingimplicationsproofs}
dataset to compare logical reasoning to the baselines.
ProofWriter is a logical reasoning dataset where each sample
contains a small rulebase (facts plus rules) and a query to evaluate as true,
false, or unknown. It is useful here because it directly tests multi-step
compositional inference, which matches the
intended behavior of the routed logic pathway.
Each input is the rulebase followed by the question
(\texttt{"\{theory\}\textbackslash nQuestion: \{question\}"}, about 120 tokens
on average), truncated from the left if needed so the question is always kept.
We use deterministic, matched sampling across all models
(10k train, 5k validation, all problem depths). Main runs are trained for 3
epochs. Validation accuracy still rose from epoch 2 to 3 in 17 of 24 runs and
fell by less than 1 pp in the rest, so longer training might raise scores
slightly.

\subsubsection{Beyond Logic -- Multi-Genre Natural Language (MNLI)}
To test performance beyond rulebase logic tasks, we evaluate
performance on Multi-Genre Natural Language (MNLI) \cite{N18-1101}.
MNLI allow us to evaluate the natural language understanding
of the model. MNLI has logical multi-class labels
(Entailment, Neutral, Contradiction),
this allows us to map to the same 3-label classification
head used for ProofWriter.

\subsection{Hyperparameter Evaluations}
The reported ablations vary gate count $G$ (8 vs 16), window axis
(inter-token vs intra-token), fusion $\alpha$ (0.01 vs 0.1, fixed vs learned),
and cross-attention (present vs removed). Logic dimension ($L=256$), window
size, and gate mode (soft) are held fixed. We validate on task validation sets
and compare against the No-Gate Control to avoid attributing gains to
unmatched model size.

\subsection{Metrics}
Training metrics are validation accuracy and validation loss.
Operator mechanism metrics are routing entropy and fusion $\alpha$. We report
routing entropy normalized by its maximum, $H/\ln(3G)$, so that runs with
different gate counts are comparable, where 1 means uniform routing.
We also report average gradient norm (Avg Grad Norm) as an optimization-health
signal. Stable gradient norms generally indicate smoother updates,
very large norms can indicate unstable training dynamics, small
gradient norms indicates arrested training.

\section{Results}

\subsection{ProofWriter}
\begin{figure*}[tbp]
\centering
\includegraphics[width=0.9\textwidth]{plots/proofwriter_accuracy.png}
\caption{ProofWriter validation accuracy after 3 epochs with the rulebase in
the input. Blue: mean $\pm$ SD over 3 seeds. Grey: individual seeds. The
y-axis starts at 0.90.}
\label{fig:proofwriter_accuracy}
\end{figure*}

Table~\ref{tab:main_results} reports final validation accuracy, validation
loss, and mechanism metrics after 3 epochs, and
Figure~\ref{fig:proofwriter_accuracy} shows the per-seed accuracies. With the
rulebase in the input,
every model reaches 0.925--0.939 accuracy, so the backbone solves most of the
task on its own. The No-Gate Control is the best model (0.939) and beats the
Baseline by $+1.4$ pp (seed-paired SE 0.2 pp). Because the No-Gate Control also
adds the fusion MLP and a normalized head, this gain should be attributed to
that whole head path, not to parallel capacity alone.

No routed variant improves on the No-Gate Control. The seed-paired differences
are $-1.3$ pp for inter-token $G=8$ with $\alpha=0.01$ (SE 0.8, per seed
$-2.5$, $+0.3$, $-1.7$), $-0.4$ pp for intra-token $G=8$ (SE 0.2), and
$-0.5$ pp for $G=16$ (SE 0.3). Changing $\alpha$ to 0.1 or learning it moves
the mean by less than one seed standard deviation. Our main hypothesis, that
inter-token routing beats the parameter-matched control, is therefore not
supported. Inter-token and intra-token composition perform the same, and
$G=16$ is no worse than $G=8$. The most tightly matched comparison is the routed
model without cross-attention against the No-Gate Control: both read the same
projected hidden state and use the same number of parameters to within 0.05\%,
so they differ only in gates versus MLP. Its seed-paired difference is
$-0.2$ pp (SE 0.35, per seed $+0.5$, $-0.7$, $-0.5$), so the gates add nothing
even with capacity matched exactly.

Gradient norms are lower for every model with the fusion head (19.6--22.9)
than for the Baseline (28.5), with no separation between routed and No-Gate
models. Routing entropy is 97--98\% of its maximum $\ln(3G)$ in every routed
run, for both $G=8$ and $G=16$, so the router spreads its weight almost
uniformly across operators (Section~\ref{sec:routing_viz}).

\subsection{Size and Transfer}

A key confound in models is added trainable parameters.
Table~\ref{tab:param_counts} lists trainable parameters and the parameters
each model uses in its forward pass. The No-Gate Control's trainable count
matches the routed variants within 0.05\%, but 71M of it is cross-attention
that the control never calls. By used parameters, the control matches the
routed model without cross-attention (within 0.05\%), and the routed models
with cross-attention use 71M more than the control. Since no routed variant
outperforms the control, including those with more used parameters, parameter
count does not hide a routing effect here. The routed model without
cross-attention performs the same as the full routed model
(Section~\ref{sec:xattn}).

\begin{table}[htbp]
\centering{
\resizebox{\linewidth}{!}{%
\begin{tabular}{|l|c|c|}
\hline
Model Variant & Trainable & Used in Forward \\
\hline\hline
Baseline & 7,504,936,963 & 7,504,936,963 \\
No-Gate Control & 7,582,192,483 & 7,510,848,355 \\
Routed (G=8) & 7,579,068,163 & 7,579,068,163 \\
Routed (G=16) & 7,579,725,059 & 7,579,725,059 \\
Routed (G=8), no cross-attn & 7,507,724,035 & 7,507,724,035 \\
\hline
\end{tabular}}}
\caption{Parameter counts. The No-Gate Control builds the 71,344,128
cross-attention parameters but never calls them. Its used count subtracts them
(the difference between the two $G=8$ routed rows). Inter- vs intra-token
composition has no effect on the count. For models with a fusion weight, the
counts are for fixed $\alpha$, which is stored as a buffer. Learned $\alpha$
adds one parameter.}
\label{tab:param_counts}
\end{table}

\subsubsection{Multi-Genre Natural Language (MNLI)}
MNLI \cite{N18-1101}. combines sentence understanding with logic multi-class labels.
We map the labels to the ProofWriter classes (important for the next test):
contradiction $\rightarrow 0$, entailment $\rightarrow 1$, neutral $\rightarrow 2$.
For comparability, evaluation uses matched sampling. Notable hyperparameters
values in the appendix.

\begin{table}[htbp]
\centering
\resizebox{\linewidth}{!}{%
\begin{tabular}{|l|c|c|}
\hline
Model & MNLI Val Acc & MNLI Val Loss \\
\hline\hline
No-Gate Control & \textbf{0.871 $\pm$ 0.006} & \textbf{0.393 $\pm$ 0.058} \\
Routed (intra, G=8) & 0.860 $\pm$ 0.033 & 0.429 $\pm$ 0.114 \\
Routed (inter, G=8) & 0.851 $\pm$ 0.048 & 0.451 $\pm$ 0.132 \\
\hline
\end{tabular}
}
\caption{Performance on MNLI, mean $\pm$ SD over 3 seeds (500 validation
examples per seed).}
\label{tab:mnli_transfer}
\end{table}

Table~\ref{tab:mnli_transfer} reports MNLI results. The No-Gate Control again
has the highest mean accuracy (0.871). The routed variants are lower
(seed-paired $-1.1$ pp intra, $-2.0$ pp inter) but not significantly
(SE 1.6 and 2.5 pp), because the routed models vary more across seeds: one
inter-token seed reaches only 0.796. At this
scale, routing gives no measurable transfer benefit on natural language
inference.

\subsection{Multi-task Training}
For multi-task training, each mini-batch is sampled from a 50/50 mixture of
ProofWriter and MNLI. This experiment tests whether routing can jointly
optimize symbolic-style and open-domain NLI supervision without degrading
either task. Both models use identical multi-task settings (listed in the
appendix), and we report per-task validation metrics after joint training.

\begin{table}[htbp]
\centering
\resizebox{\linewidth}{!}{%
\begin{tabular}{|l|c|c|c|}
\hline
Model & ProofWriter Acc & MNLI Acc & Mean \\
\hline\hline
No-Gate Control & 0.918 $\pm$ 0.017 & \textbf{0.892 $\pm$ 0.012} & 0.905 \\
Routed (inter, G=8) & \textbf{0.937 $\pm$ 0.009} & 0.882 $\pm$ 0.016 & \textbf{0.909} \\
\hline
\end{tabular}
}
\caption{50/50 mixed ProofWriter (with rulebase) and MNLI training, mean $\pm$
SD over 3 seeds.}
\label{tab:multitask_mix}
\end{table}

Table~\ref{tab:multitask_mix} reports the 50/50 multi-task evaluation. Here
the routed model is higher on ProofWriter ($+1.8$ pp seed-paired, SE 0.6) and
lower on MNLI ($-1.0$ pp, SE 0.8), for about the same mean. The ProofWriter
difference is the only comparison in this paper that favours routing, but it
should be read as exploratory: it was not a pre-specified hypothesis, it is one
of nine comparisons, and with three seeds a two-standard-error threshold is
lenient (the 95\% $t$ critical value with 2 degrees of freedom is about 4.3,
while here $t \approx 3.2$). It could indicate that routing helps keep the logical
task from being crowded out by joint training, which is worth testing with
more seeds.

\section{Discussion}

\begin{figure*}[tbp]
\centering
\includegraphics[width=0.95\textwidth]{plots/routing_heatmap.png}
\caption{Routing-weight heatmap over transformer depth for the
Routed (inter, $G=8$) model (seed 42, one validation batch, averaged over
tokens). Rows are layers and columns are the $3G=24$ routed channels (AND, OR,
NOT groups), and color is mean routing weight. Uniform routing is $1/24 \approx
0.042$, which is where most cells sit.}
\label{fig:routing_heatmap}
\end{figure*}

With the rulebase in the input, a fully fine-tuned Llama-3.1-8B already
reaches 0.925 on ProofWriter, and none of the logic-expert variants improves on
a parameter-matched control. The mechanism metrics agree with the accuracy
results. Routing entropy is near its maximum in every routed run, so the router
does not specialize operators. When $\alpha$ is learned from 0.01 it shrinks to
$0.0077 \pm 0.0014$, so training turns the logic branch down rather than up.
Together these suggest that at this scale the backbone does not need, and does
not use, the operator pathway. The inter-token versus intra-token choice, which
motivated this design, makes no measurable difference.

The only comparison that favours routing is ProofWriter accuracy under 50/50
multi-task training ($+1.8$ pp). It is exploratory and comes with an MNLI cost
($-1.0$ pp), but it is the natural place for a pre-registered follow-up with
more seeds.

Fusion $\alpha$ needs care as a diagnostic. The fusion module is cast to the
backbone's bfloat16, where the spacing between representable values near 0.1 is
about $4.9 \times 10^{-4}$. The $\alpha$ learning rate ($2 \times 10^{-4}$) is
below half of that, so every update rounds away and an $\alpha$ initialized at
0.1 never moves, even when it is nominally learned. This applies to all MNLI
and multi-task runs. Near 0.01 the spacing is small enough for $\alpha$ to move,
which is why the learned-$\alpha$ ProofWriter run shows a real trajectory.

\subsection{Cross-Attention Ablation}
\label{sec:xattn}

\begin{table}[t]
\centering
\small
\setlength{\tabcolsep}{4pt}
\resizebox{\linewidth}{!}{%
\begin{tabular}{|l|c|c|}
\hline
Model & Validation Acc & Validation Loss \\
\hline\hline
Routed (inter, G=8) & 0.927 $\pm$ 0.012 & 0.203 $\pm$ 0.021 \\
Routed (inter, G=8), no cross-attn & 0.937 $\pm$ 0.010 & 0.178 $\pm$ 0.031 \\
No-Gate Control & \textbf{0.939 $\pm$ 0.004} & \textbf{0.167 $\pm$ 0.014} \\
\hline
\end{tabular}}
\caption{ProofWriter with and without cross-attention into logic space, mean
$\pm$ SD over 3 seeds. Without cross-attention the routed stream reads a linear
projection of the backbone hidden state plus its previous state.}
\label{tab:routed_gate_only_vs_nogate}
\end{table}

Table~\ref{tab:routed_gate_only_vs_nogate} removes cross-attention from the
routed model while keeping everything else fixed. The model without
cross-attention scores $+1.1$ pp higher (seed-paired SE 1.1 pp, not
significant) with 71M fewer parameters, so cross-attention into logic space
contributes nothing measurable on this task. Without cross-attention the routed
model has the same input path as the No-Gate Control and matches its used
parameters within 0.05\%, which makes this pair the cleanest gates-versus-MLP
comparison: $-0.2$ pp (SE 0.35). The fourth cell of this design, a No-Gate
Control with cross-attention, was not run.

\subsection{Routing Visualization}
\label{sec:routing_viz}

Figure~\ref{fig:routing_heatmap} shows the routing weights of the
inter-token model. Most cells are close to the uniform value $1/24$, with
isolated cells up to about 0.17 and no consistent operator preference across
layers. A few bright cells can make the heatmap look sparse, but normalized
entropy of 0.97--0.98 shows that the router stays close to uniform.

\section{Limitations}

We evaluate one logical reasoning benchmark and one transfer task, with a single
8B backbone and three seeds per configuration. ProofWriter with the rulebase is
largely solved by the backbone alone (0.925), which leaves little headroom for
an auxiliary pathway. Harder or out-of-distribution reasoning splits, such as
compositional generalization benchmarks
\cite{keysers2020measuringcompositionalgeneralizationcomprehensive}, could
behave differently. All models train in pure bfloat16 without fp32 master
weights, which rounds away small updates (see $\alpha$ above) and may
under-train the logic modules in general. All variants share one set of
hyperparameters, and none was tuned per variant. The routing design also introduces
hyperparameters (gate count, window size, routing temperature) whose
interactions are not fully explored.

\section{Conclusion}

This architecture shifts conditional computation from routing among
independent neural experts to routing over fuzzy logical operator groups, with
matched-capacity controls and mechanism-level diagnostics. In seed-paired
comparisons on ProofWriter with the rulebase in the input, a parallel branch
with a fusion head improves on a linear-head Baseline by $+1.4$ pp in every
seed, but the gain does not depend on the logic operators: no routed variant
outperforms a parameter-matched No-Gate Control, cross-attention into logic
space adds nothing measurable, and the inter-token versus intra-token choice
does not matter. On MNLI, routed and No-Gate variants are also
indistinguishable. Near-maximal routing entropy and a shrinking learned
$\alpha$ indicate that the model barely uses or specializes the operator
pathway. Useful next steps are harder reasoning benchmarks where the backbone is not
already near ceiling, fp32 training of the logic modules, and a pre-registered
multi-task test of the one exploratory gain we observed.

{\small
\bibliographystyle{ieee}
\bibliography{egpaper_refs}
}

\appendix
\section*{Appendix}

\setcounter{figure}{0}
\setcounter{table}{0}
\setcounter{equation}{0}
\renewcommand{\thefigure}{A\arabic{figure}}
\renewcommand{\thetable}{A\arabic{table}}
\renewcommand{\theequation}{A\arabic{equation}}

\section{ProofWriter Evaluation Details}
This appendix summarizes the main ProofWriter configuration used in the
experiments.

\subsection{Settings and Hyperparameters}
\begin{itemize}
\item Batch size: 30 (train), 30 (eval)
\item Epochs: 3
\item Learning rate: $1\times10^{-5}$ (fusion $\alpha$: $2\times10^{-4}$ when learned)
\item Weight decay: 0.01
\item Logic dimension / gates / heads: 256 / 8 (16) / 4
\item Gate mode: \texttt{soft}
\item Precision: bfloat16, gradient checkpointing on
\item Input: rulebase (\texttt{theory}) then question, left-truncated, at most 250 tokens
\item Seeds: 42, 777, 123
\end{itemize}

\subsection{Dataset(s)}
\begin{itemize}
\item ProofWriter Depth: \texttt{all}: (all allowed)
\item Train samples: 10,000
\item Validation samples: 5,000
\end{itemize}

\section{MNLI Evaluation Details}
For MNLI transfer, labels are remapped into the shared
three-class order used throughout this paper:
contradiction $\rightarrow 0$, entailment $\rightarrow 1$, and
neutral $\rightarrow 2$.

\subsection{Dataset(s)}
\begin{itemize}
\item Train samples: 5,000 (\texttt{train} split)
\item Validation samples: 500 (\texttt{validation\_matched})
\end{itemize}

\subsection{Settings and Hyperparameters}
\begin{itemize}
\item Train batch size: 20
\item Eval batch size: 10
\item Epochs: 2
\item Learning rate: $1\times10^{-5}$
\item Weight decay: 0.01
\item Logic dimension / gates / heads: 256 / 8 / 4
\item Fusion $\alpha$: 0.1, nominally learned but frozen by bfloat16 rounding
\item Gate mode: \texttt{soft}
\item Seeds: 42, 777, 123
\end{itemize}


\section{Multi-task Evaluation Details}
For multi-task experiments, each batch is sampled from a mix
of ProofWriter and MNLI datasets.

\subsection{Dataset(s)}
\begin{itemize}
\item MNLI sampling ratio: 0.5
\item ProofWriter sampling ratio: 0.5
\item ProofWriter depth: \texttt{all}
\item Train/validation samples: 10,000 / 1,000
\end{itemize}

\subsection{Settings and Hyperparameters}
\begin{itemize}
\item Train batch size: 20
\item Eval batch size: 100
\item Epochs: 6
\item Learning rate: $1\times10^{-5}$
\item Weight decay: 0.01
\item Logic dimension / gates / heads: 256 / 8 / 4
\item Fusion $\alpha$: 0.1, nominally learned but frozen by bfloat16 rounding
\item Gradient checkpointing: on, seeds 42, 777, 123
\item Gate axis: \texttt{inter\_token}
\item Gate mode: \texttt{soft}
\end{itemize}

\section{Cross-Attention Ablation Details}
Table~\ref{tab:routed_gate_only_vs_nogate} uses the ProofWriter settings above
with inter-token routing, $G=8$, and fixed $\alpha=0.01$. In the ablated model
no cross-attention modules are built. At each layer the routed stream takes the
logic projection of the backbone hidden state plus the previous logic state,
the same input path as the No-Gate Control.

\section{ProofWriter Question}
This appendix section provides an example question from ProofWriter.

\subsection{Rulebase}
\begin{verbatim}
triple1: The cow is big.
triple2: The cow needs the dog.
triple3: The dog sees the rabbit.
triple4: The rabbit chases the cow.
triple5: The rabbit chases the dog.
triple6: The rabbit is big.
triple7: The rabbit sees the dog.
rule1: If the cow is blue and the cow
       needs the rabbit then the cow
       needs the dog.
rule2: If the cow chases the dog then
       the cow sees the rabbit.
rule3: If something is big then it
       chases the dog.
\end{verbatim}

\subsection{Question, Answer}
\begin{verbatim}
Question: The cow sees the rabbit?
Answer: true
\end{verbatim}

\end{document}
